% DRAFT — main_zh.tex 镜像小节(与 draft_campaign_section_en.tex 同步;T 留槽)
% 位置:插在 5.5 之后、Additional observations 之前。

\subsection{比较行为证据(战役 v2)}
\label{sec:res-campaign}

上文的机制比较使用闭式与嵌入检索对照。为在同一问题上检验行为层,
我们运行了第二轮战役(360 run;设计、两次中途混淆修复与完整数据见
补充材料 S15--S17):所有条件共享同一脚本化经历流、同一菜单、同一
LLM 决策头(温度 0)与同一决策预算——唯一差异是记忆形式。

\textbf{中性指令下的目标重拾(C)。}开卷练习获得 \texttt{oak\_planks}
与无关干扰相之后,测试提示不含任何目标("you have free choice")且
配方事实剥离。FAS 在 \textbf{9/10} 配对种子中重拾并完成被中断的目标;
全历史 LLM 基线与无记忆基线均为 \textbf{0/10}(McNemar $p{=}0.0039$;
Holm $p{=}0.0117$)。这是战役中持久关系状态唯一干净的比较级优势。
机制归因:零经历对照复现同一工作集(配方闭包先验边
\texttt{oak\_log}$\rightarrow$\texttt{oak\_planks} 在原料被观测时
再点亮链条),故效应是\emph{观测触发的先验关系结构激活}——关系
工作集显著度——而非意图生命周期机制。

\textbf{目标修订是同一机制的代价(E)。}在不可达目标(工具门槛矿石)
下,全历史基线平均 \textbf{1.0} 次失败即转向;FAS 重复 \textbf{10.5}
次,7/10 转向且偏晚。两个结果是同一枚硬币:在干扰下保持目标显著的
状态,在环境使其失效时同样保持其显著。

\textbf{消费者修复(E2/F)。}我们随后检验失败与目标状态是否只是
未被决策上下文消费。把生产事件框架的负极性编码接入战役装配
(失败$\rightarrow$实体边 $w{=}{-}0.8$,参数与生产配置一致)使重复
失败从 10.5 降至 \textbf{7.2}(单侧 Wilcoxon $p{=}0.011$)——统计
显著但仍未及基线,残余惰性属设计层(恒定目标文本、弱账本先验)。
加入统一自我图目标消费者(\texttt{current\_goal} $\rightarrow$
工作集;生产默认关)使此前逐位相同的 ON/OFF 行为发生分离:registry
清空后,consumer 开启的循环在 5/5 种子中重新派生被放弃的目标,
每 run 采集动作从 1.0 升至 \textbf{6.0}。自我模型的"行为无关"
是缺消费者伪影,而非表征性质。

\textbf{行为汇聚(I)。}双链汇聚任务(\texttt{oak\_fence} 需要分两次
经历习得的 \texttt{oak\_planks} 与 \texttt{stick})上,FAS 条件完成
0--2/10,而原始历史 LLM 基线完成 \textbf{10/10}(Holm $p{=}0.041$)。
上下文预算($k{=}16,32$)与结构增富序列化(显式
动作$\rightarrow$结果对)均为 0/10($n{=}30$),把失败定位在
\emph{序列化格式}:图谱状态丢弃了原始流保留的时序叙事
(补充材料 S17)。

\textbf{信用地平线(B)与长程稳定性(J)。}生产归因管线上的确定性
延迟探针显示,动作$\rightarrow$后果信用在 80\,s 内 10/10 形成、
160\,s 时 0/10——90\,s 归因窗即架构的信用地平线(有意的语义边界,
第 7 节)。100 周期稳定性运行(零 LLM)显示图谱有界增长
($+3.7\%$ 节点)、激活受钳制、无污染崩塌,但晋升亦为零
(重复饥饿的情节未达支持阈值)——长程巩固仍未检验。

	extbf{联合测量(T)。}在同一任务序列上同时携带两个机制(不可达目标的失败情节 + 中性自由相):事件框架编码复现其效应(失败重复 12.0$
ightarrow$7.8,单侧 Wilcoxon $p{=}0.002$);目标 consumer 在先验边路径已饱和处(全条件 10/10)无重拾增量,也无惰性惩罚(6.8 vs 7.8,n.s.)——两个消费者可共存,且组合工作集把失败动作完全挤出 top-8。抑制与显著度的行为权衡曲线可测,其前两点偏向抑制。(一处设计缺陷——库存遗留使首个重拾指标饱和——被发现后已用新世界重跑修正;两份数据均已归档。)

综合以上,本战役修正论文的主张结构:持久关系状态在一个方向上
可证明地改变行为(重拾),在另一方向上可证明地付出代价(修订),
当前瓶颈在序列化接口(汇聚),并被显式设计窗口限幅(延迟后果)。
全部反向与负结果以同等力度报告。
